\documentclass[12pt]{article}
\usepackage[left = 0.6in, right = 0.6in, top = 0.6in, bottom = 1.5in]{geometry} 
\usepackage{enumerate, pgfplots, fontawesome5, booktabs, xcolor, hyperref, amsmath,amsthm,amssymb,amsfonts, fancyhdr, color, comment, graphicx, environ, actuarialsymbol}
\pagestyle{fancy}
\setlength{\headheight}{65pt}
%%%%%%%%%%%%%%%%%%%%%%%%%%%%%%%%%%%%%%%%%%%%%
\lhead{}
\rhead{Differentiation and applications}
%%%%%%%%%%%%%%%%%%%%%%%%%%%%%%%%%%%%%%%%%%%%%
%%%%%%%%%%%%%%%%%%%%%%%%%%%%%%%%%%%%%%%%%%%%%
\newcommand{\emptygrid}[3]{ % #1: range, #2: y-label, #3: "labels" or "nolabels"
    \begin{tikzpicture}
    \begin{axis}[
        axis lines = middle,
        xlabel = $x$,
        ylabel = $#2$,
        xmin=-#1, xmax=#1,
        ymin=-#1, ymax=#1,
        grid = both,
        grid style = {line width=.1pt, draw=gray!30},
        xtick distance = 5,
        ytick distance = 5,
        minor tick num = 4,
        tick label style = {font=\tiny},
        width=8cm, height=8cm,
        axis equal image,
        % Logic to hide labels if #3 is "nolabels"
        xticklabels={ \ifx\empty#3\empty\else\ifnum\pdfstrcmp{#3}{nolabels}=0 \else \fi\fi },
        yticklabels={ \ifx\empty#3\empty\else\ifnum\pdfstrcmp{#3}{nolabels}=0 \else \fi\fi }
    ]
    \end{axis}
    \end{tikzpicture}
}
%%%%%%%%%

\begin{document}
\section{Differentiation}
\begin{enumerate}
    \item if $y = (x + \sqrt{x^2+1})^2,$ show that 7. if $y = (x + \sqrt{x^2+1})^2,$ show that $\frac{dy}{dx} = \frac{2y}{\sqrt{x^2+1}}$
        

\end{enumerate}

\section{Applications of differentiation}





   \begin{enumerate}
        \item Find the smallest angle between the tangent to the graph of $f$ at $x = 1$ and the x axis
        \item Find the acute angle between the line $y = -3x + 13$ and the tangent to $f$ at $x = \frac{1}{2}$





    \item City A is located at the point (0, 0) and city B is located at the point (10, 0). Let $$f(x) = \frac{p}{x+1}+\frac{4}{11-x}, x \in [0,10]$$ represent the pollution level between the two cities.
    \begin{enumerate}
        \item Find value(s) of $p$ such that city B is the most polluted point between the two cities.
        \item Find value(s) of $p$ such that city B is the least polluted point between the two cities.
    \end{enumerate}
    \item The concentration of insects per square metre in a gorge, $C$ is a continuous function of time $t$, given by
    $C(t) = \begin{cases} 1000 \left( \cos \left( \frac{\pi(t-8)}{2} \right) + 2 \right)^2 - 1000 & 8 \leq t \leq 16 \\ m & 0 \leq t < 8 \text{ or } 16 < t \leq 24 \end{cases}$
    
    Where $t$ is the number of hours after midnight and $m \in \mathbb{R}$ 
    \begin{enumerate}
        \item What is the value of $m$?
        \item Sketch the graph of $C$ for $t \in [0,24]$
        \item What is the minimum concentration of insects and at what value(s) of $t$ does that occur?
        \item The insect infestation is considered deadly if the concentration is more than 1250 insects per square metre. At what time after midnight does the concentration of insects first stop being deadly?
        \item During a 24 hour periiod, what is the total length of time for which the insect infestation is not deadly?
        \item Due to the uneven surface of the gorge, the time, $T$ minutes, it takes to run $x$ km is given by $T = p(q^{x} -1)$ where $p$ and $q$ are constants. 
        \begin{enumerate}
            \item If it takes 5 minutes to run the first kilometre and 12.5 mins to run the first two kilometres, find the values of $p$ and $q$
            \item Find the length of time it takes to run 4km into the gorge to retrieve a buried treasure
            \item It takes 19 minutes to dig up the treasure and you are able to run back out of the gorge in half the time it takes to reach the treasure. Show that it is possible to recover the treasure successfully and state how much time there will be to spare.
        \end{enumerate}
    \end{enumerate}
    
    \item The light intensity level $L$, in a greenhouse varies over the course of a day according the function below.
    \begin{align*}
    L(t) = \begin{cases} -\frac{1}{18}t^2 + \frac{4}{3}t - 6 & , \quad 6 \le t < 7 \cup 17 \le t < 18 \\ \frac{1}{4} \sin \left( \frac{3\pi t}{5} \right) + H & , \quad 7 \le t < 17 \\ 0 & , \quad \text{elsewhere} \end{cases}
    \end{align*}
    Where $t$ is the number of hours after midnight and $H \in \mathbb{R}^{+}$.
    \begin{enumerate}
        \item Assuming $L(t)$ is a continuous function, find the value of $H$
        \item Sketch the graph of $L$ for $t \in [0,24]$
        \item Find the maximum light intensity level and at what value(s) of $t$ does that occur?
        \item For what time interval(s) is the light intensity in the top 10\% of the intensities throughout the day?
    \end{enumerate}

\item The hobbits, on their way to the village of VCBree find themselves point A on the south side of a rectangular forest that is 2km wide, and need to get to point to the village VCBree located at point C, 20km to the east and on the other side of the forest. Point B is directly to the north of point A and on the edge of the forest as shown. There is a track running along the far side BC of the forest. They need to get from A to C. They can travel at a speed of 24km/hr on the track but only at a speed of 12km/hr through the forest. Note: The Diagram is not to scale.

\vspace{1.5em}

\begin{center}
\begin{tikzpicture}
    % Main rectangle
    \draw (0,0) rectangle (12,4);
    
    % Labels
    \node[below left] at (0,0) {A};
    \node[above left] at (0,4) {B};
    \node[above right] at (12,4) {C (Village of VCBree)};
    \node[below right] at (12,0) {D};
    \node[above] at (4,4) {E};
    
    % Diagonal line AE
    \draw[thick] (0,0) -- (4,4);
    
    % Track markings along top edge BC
    \draw (0,4.2) -- (12,4.2);
    \foreach \x in {0, 0.15, ..., 12} {
        \draw (\x, 4) -- (\x, 4.2);
    }
    
    % Dimension arrows
    \draw[<->] (-0.8,0) -- (-0.8,4) node[midway, left] {2 km};
    \draw[<->] (0,-0.8) -- (12,-0.8) node[midway, below] {20 km};
    
    % Angle theta
    \draw (1,0) arc (0:45:1);
    \node at (1.4, 0.4) {$\theta$};
    
    % Trees (using fontawesome5)
    \foreach \x/\y in {
        0.8/1.2, 1.2/2.5, 1.8/3.5, 2.5/2.2, 3/1, 
        3.5/2.8, 4.5/1.5, 5/2.5, 5.5/3.5, 6/0.8, 
        6.5/1.8, 7.5/1.2, 7.5/2.8, 8/3.5, 8.5/2, 
        9/0.8, 9.5/1.5, 10.5/1, 11/2.5, 11.5/1.5
    } {
        \node[text=green!60!black, scale=1.5] at (\x,\y) {\faTree};
    }
\end{tikzpicture}
\end{center}

\vspace{1em}

\begin{enumerate}
    \item Calculate the time taken to travel directly from A through the forest to B and then from B to C using the track.
    \item Find the distance the hobbits travel inside the forest from $A$ to $E$ as a function of $\theta $
    \item Explain why the possible values of $\theta $, correct to the nearest degree, are $\theta  \in [6^{o}, 90^{o}]$
    \item Find the time taken to travel from $A$ to $C$ through the forest, as a function of $\theta $
    \item Find the value of $\theta $ that minimises the time taken to travel from $A$ to $C$ through the forest and calculate the minimum time taken.
\end{enumerate}
\end{enumerate}
\end{document}